% discussion.tex
% Discussion: scope, caveats, and extensions.

In this study, we composed synthetic per-asset growth-rate paths
under a single-index model (SIM) by drawing each asset's innovation
from an independent JumpHMM generator fit to the asset's own
full growth-rate history, then rescaling the draw by a closed-form scalar
before inserting it into the SIM composition. We evaluated six
composers on a calibrated universe of $424$ US equities and
exchange-traded funds against the real SPY market path: the naive
composition, a Gaussian-residual SIM baseline, the hybrid
construction proposed here, and three residual-fit alternatives
(JumpHMM-on-residuals, Politis--Romano block
bootstrap~\cite{politisRomano1994}, and a per-ticker GARCH(1,1)-$t$
fit~\cite{bollerslev1986}). The one-line variance correction,
together with the sign-preserving clip rule and the $R^2$-preserving
branch, closes a gap between two previously unsatisfying options for
building synthetic asset paths with both marginal structure and
factor structure: naive composition gives both but distorts the
marginal by $\beta^2 \sigma_m^2$, while Gaussian-residual SIM gives
the factor exactly but strips the marginal. The hybrid construction
sat between the Gaussian-residual and naive baselines: it recovered
$\beta$ to the same precision as either, held the generator marginal
to its target variance in every ticker, and preserved the generator's
tails at a level indistinguishable from the naive composition. The
three residual-fit composers produced higher Kolmogorov--Smirnov
pass rates than hybrid on marginal body fidelity ($67$ to $76\%$
versus $50\%$), but matched hybrid on the $99\%$ Value-at-Risk
coverage test to within statistical tolerance
(Table~\ref{tab:var}); the body-fidelity gap did not translate to a
tail-accuracy gap. Hybrid's contribution is therefore scoped: it is
the best composition rule when the per-asset generator is specified
in full-return units (a common setting, since per-asset generators
are often calibrated for distributional interpretability, cached for
other pipelines, or chosen to retain regime-switching or jump
structure), and it delivers $99\%$ VaR accuracy on par with a
dedicated residual generator without requiring a separate fitting
stage. The gap was large
enough to matter in practice: at the high-$\beta$ end of the
universe, naive composition inflated the marginal variance by over
$70\%$ relative to the generator target (Fig.~\ref{fig:preservation}). Because branch selection is
driven entirely by the calibrated $R^2$ and each asset carries a
persisted construction flag, the procedure runs uniformly across an
arbitrary ticker universe without external metadata and leaves an
audit trail that downstream consumers can inspect. No part of the
construction is new in a deep sense; the contribution is a minimal,
practitioner-ready recipe with an auditable branch flag and a
principled treatment of the degenerate tracker limit, usable as a
drop-in replacement for naive composition in workflows where both
tail shape and market dependence matter.

The evaluation isolated the composition rule by fixing the market
path at the real SPY history and using the
same per-ticker innovation draw across all three composers. This design
yielded a paired comparison of composition rules at the per-ticker level
but set aside two questions that production users will need to address
separately. The first concerns cross-sectional dependence: we skipped
the optional rank-reorder step in Algorithm~\ref{alg:hybrid} so that
per-ticker metrics reflected the composition rule alone. When the
reorder is applied, the variance algebra is unchanged
(rank-reordering preserves the marginal variance of $\eps_i$), so the
marginal results here carry over; cross-sectional fidelity under
different copula specifications is studied in the companion
\texttt{JumpHMM.jl} paper~\cite{alswaidanVarner2026jdiq}. The second
concerns market-path uncertainty: the metrics reported here used the
real SPY path, and under a simulated market path the composed
marginal picks up additional variance from the sampling of
$g_m^{(r)}$, which for high-$R^2$ tickers becomes a noticeable
effect. In production, the two deferred questions compose:
a reorder-plus-simulated-market pipeline would layer cross-sectional
structure on top of per-ticker draws built against sampled market
paths, with the variance correction of Eq.~\eqref{eq:scale} still
providing the per-asset variance target on each realized
$g_m^{(r)}$. A separate scope boundary concerns the downstream use
case: the main metrics reported here characterize marginal and
factor fidelity of the composition rule on per-ticker targets. To
test whether that fidelity translates to a consequential downstream
number, we also report a per-ticker one-day $95\%$ and $99\%$
Value-at-Risk backtest of the composed paths against the real
2014--2024 growth-rate histories, with exceedance counts assessed by
Kupiec's unconditional coverage test; the per-ticker framing keeps
this validation inside the paper's scope, while portfolio-level VaR
requires the cross-sectional copula deferred to the companion
paper~\cite{alswaidanVarner2026jdiq}. High-turnover portfolio
strategies remain a distinct scope boundary: the temporal structure
of the residual draw becomes load-bearing in a way that
contemporaneous composition does not address.

Clipping did not trigger anywhere in the US-equity universe used here,
because the JumpHMM marginals carried enough variance that
$\rho_i = \beta_i^2 \sigma_m^2 / \sigma_{\gen,i}^2$ stayed below $0.90$
even for the highest-$\beta$ ticker. Clipping becomes relevant in two
settings we did not evaluate: scenario-driven markets, where the
synthetic $g_m$ is constructed to be much more volatile than the
calibration market (e.g., stress tests), and universes that mix
low-volatility single stocks with extreme betas (e.g., leveraged ETFs).
In both settings we recommend reporting the original and effective
$\beta_i$ alongside the output, as in Algorithm~\ref{alg:hybrid}, so
that downstream consumers can audit the attenuation. A related
attenuation appeared at the other end of the $\beta$ spectrum, even
when clipping did not trigger: the hybrid KS pass rate fell from
$67\%$ at low $\beta$ to $44\%$ at high $\beta$
(Table~\ref{tab:beta_bucket}), a mild attenuation that followed
directly from the composition algebra. As
$\beta_i^2 \sigma_m^2$ grows toward the idiosyncratic floor
$(1 - f)\sigma_{\gen,i}^2$, the market-driven component carries an
increasing share of the marginal, and that component is only as
heavy-tailed as $g_m$. The market's own heavy tails slow the attenuation
(SPY growth rates are themselves heavy-tailed after JumpHMM fitting),
but the effect is visible. Users who need exact marginal fidelity at
high $\beta$ can recover it by lowering the floor $f$ (at the cost of
less conservative clipping behavior) or by pre-scaling the generator
variance so that $\rho_i$ stays small.

The variance correction $s_i^2 = 1 - \rho_i$ relies on
$\mathrm{Cov}(\teps_i, g_m) \approx 0$. The JumpHMM marginals are fit
in isolation and sampled independently of the market path, so this
holds by construction; if the generator is already correlated with
the market (for example, a bootstrap from joint history), the
empirical covariance should be subtracted from $\beta_i$ before
applying the correction, or $s_i^2$ should be estimated directly from
the empirical variance of the composed series. The same warning
applies to the $R^2$-preserving branch, whose variance identity
assumes $\eps$ is uncorrelated with $g_m$ after any reorder. The
construction also injects a contemporaneous market factor only:
$g_i(t)$ depends on $g_m(t)$ and on $\teps_i(t)$ but not on
$g_m(t-1), g_m(t-2), \dots$, so lagged market dependence (lead-lag
effects, market-driven volatility clustering, asymmetric responses to
drawdowns) is not reproduced by this recipe. The natural extension is
a lagged-factor SIM,
$g_i(t) = \alpha_i + \sum_{k=0}^K \beta_{i,k}\, g_m(t - k) + \eps_i(t)$;
the variance correction generalizes to that form by replacing
$\beta_i^2 \sigma_m^2$ with the variance of the full market-factor
component, though the empirical value of the extra parameters is
beyond the scope here. Lagged idiosyncratic dependence is likewise
absent: the JumpHMM residuals $\teps_i$ are drawn independently
across time, so any residual autocorrelation present in the real
data is discarded by the composition. This matters most for
workflows that evaluate portfolios rebalancing at high frequency.
With iid residuals, a rebalancing rule collects the full
diversification return of \cite{boothFama1992}, a compounding
premium that scales with residual variance and rebalancing frequency
and is largely absent in real markets, where daily residual
autocorrelation is small but positive \cite{boucheyGunthorpSutherland2012}.
Two generator alternatives in Table~\ref{tab:aggregate} preserve
temporal structure directly: the block-bootstrap composer resamples
blocks of the empirical OLS residual series~\cite{politisRomano1994},
inheriting whatever autocorrelation is present in the data, and the
GARCH(1,1)-$t$ composer fits a conditional-variance model per
ticker~\cite{bollerslev1986}. Both recover marginal fidelity at the
cost of committing the pipeline to a specific residual generator,
whereas the variance correction of Eq.~\eqref{eq:scale} is
generator-agnostic and accepts a full-return JumpHMM marginal (with
its jump and regime structure) as input. Users evaluating
high-turnover strategies on JumpHMM-corrected paths should benchmark
absolute return levels against a static-allocation baseline on the
same paths, or adopt the bootstrap or GARCH alternatives when
residual autocorrelation is the load-bearing structure for the
downstream task. Two further extensions follow naturally. Multi-factor SIM replaces
the single market factor with $K$ factors, requiring $\rho_i$ to
become the sum of factor contributions and the loading vector to be
clipped rather than a scalar, though the branch logic is unchanged.
Per-ticker tuning of the JumpHMM jump parameters $(\epsilon, \lambda)$
tightens each marginal's fit; we used defaults to keep fitting
tractable across the universe, but tuning as reported in the
companion paper~\cite{alswaidanVarner2026jdiq} would push the hybrid
KS pass rates upward. A
stress-test harness ties these together: using the clipping rule to
enforce a minimum idiosyncratic share in constructed market scenarios
guarantees that per-asset marginals remain heavy-tailed even under
synthetic markets much more volatile than history, which is precisely
the regime where naive SIM composition breaks down most severely.
